%% Abstract page IGDTUW B.Tech Minor Project Report
\begin{spacing}{1.4}
\noindent Organized financial fraud in digital credit, e-commerce, and decentralized peer-to-peer (P2P) lending has rapidly evolved from uncoordinated individual defaults into sophisticated, multi-party syndicates. Modern fraud rings systematically evade detection through deliberate topological camouflage—establishing benign connections to legitimate borrowers, spoofing IP/device fingerprints, and dispersing transactions across multi-hop subgraphs. While Graph Neural Networks (GNNs) have emerged as the state-of-the-art paradigm for learning over relational topologies, existing anti-fraud architectures (such as CARE-GNN, BWGNN, and GAGA) suffer from three critical bottlenecks: high sensitivity to class imbalance, computational inefficiency of reinforcement-learning-based neighbor filtering, and an absence of persistent structural memory across message-passing iterations.\par

\noindent To overcome these limitations, this project introduces \textbf{ACO-GNN}, a bio-inspired hybrid framework that integrates \textbf{Ant Colony Optimization (ACO) stigmergic principles} directly into the graph message-passing layers. By deploying virtual ant agents to traverse multi-relational graphs, the system establishes continuous artificial pheromone trails ($\tau_{ij}$) along high-risk paths. These persistent pheromone weights are blended with feature attention scores via an adaptive gating parameter ($\lambda$), dynamically scaling aggregation weights to neutralize camouflage edges while accentuating multi-hop fraud rings. The framework is paired with Focal Loss to counteract extreme class imbalance (5\%--15\% fraud prevalence) and validation-calibrated decision thresholds.\par

\noindent Extensive empirical benchmarking across real-world graphs (YelpChi with 8.05M edges, Amazon with 9.55M edges) and a controlled Synthetic Lending benchmark validates the efficacy of our model. On YelpChi, ACO-GNN outperforms the reproduced published baseline of CARE-GNN, elevating Macro-F1 from $0.6347$ to $\mathbf{0.7022}$ ($+10.6\%$ relative gain) and AUC-ROC from $0.7734$ to $\mathbf{0.8307}$. Furthermore, under adversarial camouflage stress testing (0\%--50\% synthetic noise injection), ACO-GNN maintains robust fraud detection ($F_1 > 0.76$) whereas vanilla GAT degrades drastically to $0.13$, proving that stigmergic pheromone evaporation acts as an autonomous structural filter. Finally, a progressive dynamic streaming simulation establishes that stigmergic memory adapts robustly as transaction topologies expand, confirming that pheromone-guided message passing offers a powerful new inductive paradigm for secure financial intelligence.\par
\end{spacing}

\vfill
\noindent\textbf{Keywords:} Financial Fraud Detection, Graph Neural Networks, Ant Colony Optimization, Stigmergic Pheromone Memory, Camouflage Robustness, Multi-Hop Fraud Rings.
